%
\begin{isabellebody}%
\setisabellecontext{Week{\isacharunderscore}{\kern0pt}{\isadigit{2}}}%
%
\isadelimtheory
%
\endisadelimtheory
%
\isatagtheory
%
\endisatagtheory
{\isafoldtheory}%
%
\isadelimtheory
%
\endisadelimtheory
%
\begin{isamarkuptext}%
\ExerciseSheet{7}{}%
\end{isamarkuptext}\isamarkuptrue%
%
\begin{isamarkuptext}%
\Exercise{Fold function}
  The fold function is a very generic function, that can be used to express 
  multiple other interesting functions over lists.%
\end{isamarkuptext}\isamarkuptrue%
%
\begin{isamarkuptext}%
Have a look at Isabelle/HOL's standard function \isa{fold}, defined as follows:%
\end{isamarkuptext}\isamarkuptrue%
\isacommand{fun}\isamarkupfalse%
\ fold{\isacharcolon}{\kern0pt}{\isacharcolon}{\kern0pt}{\isachardoublequoteopen}{\isacharparenleft}{\kern0pt}{\isacharprime}{\kern0pt}a\ {\isasymRightarrow}\ {\isacharprime}{\kern0pt}b\ {\isasymRightarrow}\ {\isacharprime}{\kern0pt}b{\isacharparenright}{\kern0pt}\ {\isasymRightarrow}\ {\isacharprime}{\kern0pt}a\ list\ {\isasymRightarrow}\ {\isacharprime}{\kern0pt}b\ {\isasymRightarrow}\ {\isacharprime}{\kern0pt}b{\isachardoublequoteclose}\ \isakeyword{where}\isanewline
{\isachardoublequoteopen}fold\ f\ {\isacharbrackleft}{\kern0pt}{\isacharbrackright}{\kern0pt}\ acc\ {\isacharequal}{\kern0pt}\ acc{\isachardoublequoteclose}\isanewline
{\isacharbar}{\kern0pt}\ {\isachardoublequoteopen}fold\ f\ {\isacharparenleft}{\kern0pt}x\ {\isacharhash}{\kern0pt}xs{\isacharparenright}{\kern0pt}\ acc\ {\isacharequal}{\kern0pt}\ fold\ f\ xs\ {\isacharparenleft}{\kern0pt}f\ x\ acc{\isacharparenright}{\kern0pt}{\isachardoublequoteclose}%
\begin{isamarkuptext}%
Write a function to compute the sum of the elements of a list. 
  Define two versions, one direct recursive definition, and one using fold.
  Show that both are equal.

  Hint: use automation!%
\end{isamarkuptext}\isamarkuptrue%
\isacommand{fun}\isamarkupfalse%
\ list{\isacharunderscore}{\kern0pt}sum\ {\isacharcolon}{\kern0pt}{\isacharcolon}{\kern0pt}\ {\isachardoublequoteopen}nat\ list\ {\isasymRightarrow}\ nat{\isachardoublequoteclose}\ \isanewline
\ \ \isanewline
\isacommand{definition}\isamarkupfalse%
\ list{\isacharunderscore}{\kern0pt}sum{\isacharprime}{\kern0pt}\ {\isacharcolon}{\kern0pt}{\isacharcolon}{\kern0pt}\ {\isachardoublequoteopen}nat\ list\ {\isasymRightarrow}\ nat{\isachardoublequoteclose}\isanewline
\ \ %
\isadelimproof
%
\endisadelimproof
%
\isatagproof
%
\endisatagproof
{\isafoldproof}%
%
\isadelimproof
\isanewline
%
\endisadelimproof
\isacommand{lemma}\isamarkupfalse%
\ {\isachardoublequoteopen}list{\isacharunderscore}{\kern0pt}sum\ xs\ {\isacharequal}{\kern0pt}\ list{\isacharunderscore}{\kern0pt}sum{\isacharprime}{\kern0pt}\ xs{\isachardoublequoteclose}\isanewline
\ \ %
\isadelimproof
%
\endisadelimproof
%
\isatagproof
%
\endisatagproof
{\isafoldproof}%
%
\isadelimproof
%
\endisadelimproof
%
\begin{isamarkuptext}%
\Exercise{Tail recursive \isa{reverse}}
In Isabelle, there is the reverse function \isa{rev}, which, given a list, computes another list
with the same elements but in reverse order.
Take a loot into its definition:%
\end{isamarkuptext}\isamarkuptrue%
\isacommand{thm}\isamarkupfalse%
\ rev{\isachardot}{\kern0pt}simps%
\begin{isamarkuptext}%
Recall tail recursion: a function is tail recursive if in all recursive calls it does not call 
  anything after itself\footnote{See \url{https://en.wikipedia.org/wiki/Tail_call.} for a quick review.}.
  For instance, the implementaion of \isa{rev} in Isabelle is not tail recursive, because it calls
  \isa{{\isacharparenleft}{\kern0pt}{\isacharat}{\kern0pt}{\isacharparenright}{\kern0pt}} after it calls \isa{rev} in the second equation.
  
  Write an alternative implementation of the reverse function that is tail recursive.

  Hint: recall that you can most of the times derive a tail recursive function using an accumulator
        argument.%
\end{isamarkuptext}\isamarkuptrue%
\isanewline
\isanewline
\isanewline
\isacommand{fun}\isamarkupfalse%
\ rev{\isacharunderscore}{\kern0pt}tr{\isacharcolon}{\kern0pt}{\isacharcolon}{\kern0pt}{\isachardoublequoteopen}{\isacharprime}{\kern0pt}a\ list\ {\isasymRightarrow}\ {\isacharprime}{\kern0pt}a\ list{\isachardoublequoteclose}\ \isakeyword{where}%
\begin{isamarkuptext}%
Show that the two implementations are equivalent.%
\end{isamarkuptext}\isamarkuptrue%
%
\isadelimproof
%
\endisadelimproof
%
\isatagproof
%
\endisatagproof
{\isafoldproof}%
%
\isadelimproof
\ \isanewline
%
\endisadelimproof
\isacommand{lemma}\isamarkupfalse%
\ {\isachardoublequoteopen}rev{\isacharunderscore}{\kern0pt}tr\ xs\ {\isacharequal}{\kern0pt}\ rev\ xs{\isachardoublequoteclose}%
\isadelimproof
%
\endisadelimproof
%
\isatagproof
%
\endisatagproof
{\isafoldproof}%
%
\isadelimproof
%
\endisadelimproof
%
\isadelimtheory
%
\endisadelimtheory
%
\isatagtheory
%
\endisatagtheory
{\isafoldtheory}%
%
\isadelimtheory
%
\endisadelimtheory
%
\end{isabellebody}%
\endinput
%:%file=Week_2.tex%:%
%:%19=6%:%
%:%23=8%:%
%:%24=9%:%
%:%25=10%:%
%:%29=13%:%
%:%31=15%:%
%:%32=15%:%
%:%33=16%:%
%:%34=17%:%
%:%36=20%:%
%:%37=21%:%
%:%38=22%:%
%:%39=23%:%
%:%40=24%:%
%:%42=29%:%
%:%43=29%:%
%:%44=30%:%
%:%44=35%:%
%:%45=36%:%
%:%46=36%:%
%:%47=37%:%
%:%58=45%:%
%:%61=46%:%
%:%62=46%:%
%:%63=47%:%
%:%78=51%:%
%:%79=52%:%
%:%80=53%:%
%:%81=54%:%
%:%83=57%:%
%:%84=57%:%
%:%86=60%:%
%:%87=61%:%
%:%88=62%:%
%:%89=63%:%
%:%90=64%:%
%:%91=65%:%
%:%92=66%:%
%:%93=67%:%
%:%94=68%:%
%:%96=77%:%
%:%97=78%:%
%:%98=79%:%
%:%99=80%:%
%:%100=80%:%
%:%102=85%:%
%:%115=108%:%
%:%118=109%:%
%:%119=109%:%
